\documentclass{article}
\usepackage{amsmath, amssymb, cancel, mathtools, bm}
\usepackage[left=.5in, right=.5in, top=1in, bottom=1in]{geometry}
\usepackage[most]{tcolorbox}
\usepackage[skip=1em,indent=0pt]{parskip}
\setlength{\parindent}{0pt}
\newcommand{\statvec}[1]{\underset{\sim}{\bm{#1}}} % Vector symbol (tilde under vector)
\DeclareMathOperator{\EX}{\mathbb{E}} % Expected Value symbol
\newcommand{\indep}{\perp\!\!\!\!\perp} % Independence symbol
\newcommand{\real}{\mathbb{R}} % Simplified 'Reals' indicator
\newcommand{\pto}{\overset{P}{\to}}
\newcommand{\asto}{\overset{a.s.}{\to}}
\newcommand{\rthto}{\overset{L^r}{\to}}
\newcommand{\dto}{\overset{D}{\to}}
\newcommand{\ext}{\EX_{\theta}}
% Force all aggregate symbols to always put values above/below
\let\Oldint=\int
\let\Oldsum=\sum
\let\Oldprod=\prod
\let\Oldbigcup=\bigcup
\let\Oldbigcap=\bigcap
\let\Oldlim=\lim
\renewcommand{\int}{\Oldint\limits} 
\renewcommand{\sum}{\Oldsum\limits} 
\renewcommand{\prod}{\Oldprod\limits} 
\renewcommand{\bigcup}{\Oldbigcup\limits} 
\renewcommand{\bigcap}{\Oldbigcap\limits} 
\renewcommand{\lim}{\Oldlim\limits} 
\newcommand{\infint}{\int_{-\infty}^{\infty}}
\newcommand{\iiddist}{\overset{\mathrm{iid}}{\sim}}
\newcommand{\norm}[1]{\left\lVert#1\right\rVert}
\newcommand{\boldred}[1]{\textbf{\textcolor{red}{#1}}}


\begin{document}

\section{}

$X_i \iiddist Gamma(\theta, k), k$ is known

I'm assuming this is for an $H_0: \theta = \theta_0, H_1: \theta \neq \theta_0$ test. 

(a) Find asymptotic level $1-\alpha$ confidence inerval for $\theta$ by inverting acceptance region of an asymptotic size $\alpha$ LRT

$L(\theta|x,k) = \prod \frac{k^\theta}{\Gamma(\theta)}x^{\theta - 1}e^{-kx} = \frac{k^{n\theta}}{\Gamma(\theta)^n}x^{n\theta - n}e^{-k \sum x}$

Under null, $\hat{\theta}_R = \theta_0$

LRT: $\lambda(x) = \frac{L(\hat{\theta}_R|x,k)}{L(\hat{\theta|x,k})} = $

What is $\hat{\theta}$??? No closed form for the MLE for $\theta$.

\boldred{On this part we can develop an acceptance region and inverted test by just assuming we know the MLE of $\theta$. However, to actually find an interval we would need to examine the LRT to see how it behaves}

(b)

$\EX[X] = \alpha \theta$. $M_1 = \bar{X}$

$\theta_{MoM}: \alpha \theta = \bar{X} \implies \theta_{MoM} = \frac{\bar{x}}{k}$

According to CLT, $\sqrt{n}(\bar{X} - k\theta) \sim N(0, \theta k^2)$. 

$\implies$ by Cont. Map. Theorem: $\sqrt{n}(\hat{\theta} - \theta) \dto N(0, \theta)$

Variance is $\theta$

(c)

We need $g'(\theta)^2 = \frac{1}{\theta} \implies g'(\theta) = \sqrt{\frac{1}{\theta}}$. Thus, $g(\theta) = 2\sqrt{\theta}$

(d)

$[g^{-1}(g(\hat{\theta}) - \frac{1}{\sqrt{n}}Z_{1-\alpha/2}), g^{-1}(g(\hat{\theta}) + \frac{1}{\sqrt{n}}Z_{1-\alpha/2})]$

$P(-Z_{\alpha/2} < \sqrt{n}(2\sqrt{\hat{\theta}} - 2\sqrt{\theta}) < Z_{\alpha/2}) \approx 1-\alpha$

$P(\frac{-Z_{\alpha/2}}{2\sqrt{n}} \le \sqrt{\hat{\theta}} - \sqrt{\theta} \le \frac{Z_{\alpha/2}}{2\sqrt{n}})$

$(\sqrt{\hat{\theta}} - \frac{Z_{\alpha/2}}{2\sqrt{n}})^2 \le \theta \le (\sqrt{\hat{\theta}} + \frac{Z_{\alpha/2}}{2\sqrt{n}})^2$

\section{}

$f(x|k,\theta) = \frac{\theta k^{\theta}}{x^{\theta+1}} = \theta (\frac{k}{x})^{\theta}\frac{1}{x}$

$L(k,\theta|x) = \prod \frac{\theta k^{\theta}}{x^{\theta+1}} = \frac{\theta^n k^{n\theta}}{\prod x^{\theta + 1}}$

(a) $\theta$ is known, $H_0: k = 1, H_1: k = 2$. Find most powerful size $\alpha$ test.

We can ignore the case when $x<1$ because under both hypothesis this is not possible. 

LRT: $\frac{L(k=1|x,\theta)}{L(k=2|x,\theta)} = \frac{\frac{\theta^n 1^{n\theta}}{\prod x^{\theta + 1}} I(x > 2)}{\frac{\theta^n 2^{n\theta}}{\prod x^{\theta + 1}}I(x>1)} = (\frac{1}{2})^{n\theta}I(X_{(1)} \ge 2)$

$\lambda(x) = \begin{cases} 0, & X_{(1)} \le 2 \\ (\frac{1}{2})^{n\theta}, & X_{(1)} > 2 \end{cases}$

$\phi(x) = \begin{cases} 1, & \lambda(x) > k \\ \gamma, &\lambda(x) = k \\ 0, & \lambda(x) < k \end{cases}$

$\alpha = P(\lambda(x) > k) + \gamma P(\lambda(x) = k)$

$P(\lambda(x) > k) = P(X_{(1)} > 2)$. We go from the test function to the LRT. LRT can only be greater than some k when the LRT provides a value, which only happens when the min is greater than 2. 

If we pick k = 0, the second term is also active. If $k = 2^{n\theta}$ then only the second term is active.

$P(X_{(1)} > 2) = P(x_1 > 2)^n = [\int_2^{\infty} \theta x^{-\theta - 1} dx]^n = [-x^{-\theta}|_2^{\infty}]^n = 2^{-n\theta}$

We cannot get a set $\alpha$ test using this specific test, so we will need the $\gamma$ term meaning we must let k be one of the boundary values. 

(b) k is known. $H_0: \theta \le \theta_0, H_1: \theta > \theta_0$

Need UMP size $\alpha$ test. 

Check for MLR: $T = \sum ln(\frac{X_i}{k})$

$\theta_2 > \theta_1: \frac{L(\theta_2)}{L(\theta_1)} = (\frac{\theta_2}{\theta_1})^n e^{-T(\theta_2 - \theta_1)}$. This is monotonically decreasing in $\theta_2$.

Karlin-Rubin: $\phi(x) = \begin{cases} 1, & T < c \\ 0, & T > c \end{cases}$. Choose c such that $P_{\theta_0}(T < c) = \alpha$.

\section{}

$f(x|\lambda) = \frac{2x}{\lambda^2}e^{-\frac{x^2}{\lambda^2}}; x,\lambda >0$

$L(\lambda|x) = \prod 2x \frac{1}{\lambda^{2n}} e^{-\frac{1}{\lambda^2} \sum x^2}$

(a) T = $\sum X_i^2$. Show it's complete and sufficient. 

$f(x|\lambda) = [\prod 2x] [\frac{1}{\lambda^{2n}}e^{-\frac{1}{\lambda^2}T}] = h(x)g(T, \lambda)$ thus by factorization theorem this is sufficient. Because k is known, this is actually part of the exponential family (also can be shown by writing it as $h(x)c(\lambda)e^{w(\lambda)t(x)})$. Because this is part of the exponential family, and the support contains an open set, this is a complete statistic as well. 

Still not too sure how to do pivots. This is a scale family, so the pivot would be something like $\lambda x$, but I'm not sure if that's correct or how to proceed. 

$T(X) =\sum x^2: Y = X^2 \implies Y \sim Exp(\lambda^2) \implies \sum x^2 \sim Gamma(n 1/\lambda^2)$

Scale family, need to adjust by a multiple of $\lambda^2$.

$Q = \frac{\sum X^2}{\lambda^2} \sim Gamma(n, 1)$

$P(a \le Q \le b) = 1-\alpha$

$P(\sqrt{\frac{\sum X^2}{b}} \le \lambda \le \sqrt{\frac{\sum X^2}{a}})$

Those two equations give us our system to solve and allow us to find the minimum width (Lagrangians)

(b) $H_0: \lambda = \lambda_0, H_1: \lambda \neq \lambda_0$

$l(\lambda|x) = \sum ln(2x) - 2n ln(\lambda) - \frac{\sum x^2}{\lambda^2}$

$\frac{\partial l}{\partial \lambda} = -\frac{2n}{\lambda} + \frac{2\sum x^2}{\lambda^3} = 0 \implies 2n\lambda^2 = 2\sum x^2 \implies \hat{\lambda} = \sqrt{\frac{\sum x^2}{n}} = \sqrt{\frac{T}{n}}$

LRT: $\frac{L(\lambda_0|x)}{L(\hat{\lambda|x})} = \frac{\prod 2x \frac{1}{\lambda_0^{2n}} e^{-\frac{1}{\lambda_0^2} \sum x^2}}{\prod 2x \frac{1}{\sqrt{T/n}^{2n}} e^{-\frac{1}{\sqrt{T/n}^2} \sum x^2}} = (\frac{\sqrt{T/n}}{\lambda_0})^{2n} e^{\sum x^2 (\frac{n}{T} - \frac{1}{\lambda_0^2})}$

This is monotonically increasing in $T$, thus the rejection region of $a \le \lambda(x) \le b$ is equivalent to saying $a^* \le T(x) \le b^*$

\end{document}